%!TEX program = lualatex
\documentclass[ngerman,headheight=70pt]{scrartcl}
\usepackage[ngerman]{babel}
\usepackage{blindtext}
\usepackage[utf8]{luainputenc}
\usepackage{scrlayer-scrpage}
\usepackage[T1]{fontenc}
\usepackage{fontspec}
\usepackage{multicol}
\usepackage{enumitem}
\usepackage{epigraph}
\usepackage[german=quotes,autostyle=true,strict=true]{csquotes}
\usepackage[bottom=2cm,footskip=8mm]{geometry}
\setmainfont{Source Sans Pro}
\MakeOuterQuote{"}

\setlength{\parindent}{0pt}
\setlength{\parskip}{1pt}
\newenvironment{myitemize}{\begin{itemize}\itemsep -2pt}{\end{itemize}}


\pagestyle{plain.scrheadings}
\thispagestyle{scrheadings}
\lohead{\vspace*{-1.5cm}\includegraphics[scale=0.4]{../images/stupa_logo_rot.pdf}}
\rohead{\\Präsidium des\\Studierendenparlaments}

\lofoot[Universität Hamburg · Präsidium des Studierendenparlaments\\\hrule Von-Melle-Park 5 · D-20146 Hamburg]{Universität Hamburg · Präsidium des Studierendenparlaments\\\hrule Von-Melle-Park 5 · D-20146 Hamburg}
\cofoot[]{}
\cohead[\thepage]{}

\renewcommand{\thesubsection}{TOP \arabic{subsection}}
\renewcommand{\thesubsubsection}{\arabic{subsubsection}.}
\setcounter{subsection}{-1}

\begin{document}
    UHH · StuPa-Präsidium · Von-Melle-Park 5 · D-20146 Hamburg

    \section*{Beschluss des Studierendenparlamentes vom 26. Januar 2017}
    \subsubsection*{Wissenschaftsrat bleibt hinterm Mars.\\
    Jetzt erst recht: Für eine bedarfsgemäße Finanzierung der
    Universität zum Allgemeinwohl.}

    Am 23. Januar 2017 veröffentlichte der Wissenschaftsrat seine "Empfehlung
    zur Weiterentwicklung der Geistes- und Sozialwissenschaften sowie der
    Gesamtstrategie der Universität Hamburg". Dieses Schriftstück von gut 150
    Seiten will das Ergebnis einer Begehung der Uni Hamburg im letzten halben
    Jahr sein. Es ist aber vielmehr ein weiterer Beweis dafür, dass der
    Wissenschaftsrat keine geeignete Einrichtung für eine demokratische und
    gesellschaftsverantwortliche Wissenschaftsentwicklung ist. Statt einer
    permanenten Evaluierungsshow mit Bewertungsdruck muss die Wissenschaft
    bedarfsgemäß finanziert werden, damit sie umfassend zu einer humanen
    Entwicklung der Gesellschaft beitragen kann.

    Der Wissenschaftsrat ist ein Beratungsgremium aus exklusiv ausgewählten
    Wissenschaftler*innen, Vertreter*innen aus Bund und Ländern sowie "Personen
    öffentlichen Lebens" sprich Industrie- und Großunternehmen-Vertreter*innen,
    das im Jahr 1957 für die Koordination von Großforschungsprojekten gegründet
    wurde. Von Beginn an hat er sich als Instrument für marktkonforme Wissenschaft
    hervorgetan: Seit seiner Gründungszeit, also in der Zeit der Remilitarisierung
    und NATO-Integration Westdeutschlands, bildete die Förderung von rüstungsrelevanter
    Forschung und nuklearer Technologien eine Orientierung für den Wissenschaftsrat.
    Dem gewachsenen gesellschaftlichen Bedarf an Hochschulbildung versuchte er im
    Jahr 1966 mit dem Vorschlag, dass Studium in Masse (berufsbefähigenden
    Studienabschluss) und Elite (forschungsorientiertes Studium) zu trennen, zu
    begegnen. Der breite Protest aus der Studierendenschaft in den "68er"-Jahren
    drängte diese Vorschläge zurück und setzte stattdessen einen tendenziell
    kritischen Gesellschaftsbezug in der Wissenschaft durch. Infolgedessen wurde
    der Wissenschaftsrat auch kein Planungs- sondern "nur" ein Beratungsorgan.

    Vor diesem Hintergrund kritisiert das Studierendenparlament der Uni Hamburg
    die Entscheidung des Hamburgischen Senats, den Wissenschaftsrat zur Beurteilung
    und Beratung der Entwicklung der Uni Hamburg hinzugezogen zu haben und fordert
    alle dazu auf, die Empfehlungen sehr kritisch zu lesen.

    Die Mitglieder der Uni Hamburg engagieren sich seit Jahren in der Studienreform
    zur Bildung mündiger Persönlichkeiten statt Bachelor-Master-Terror, für den
    Ausbau der demokratischen Hochschule statt Top-Down-Strukturen mit
    Hochschulrat und für die bedarfsgemäße Grundfinanzierung von Wissenschaft zum
    Allgemeinwohl statt Konkurrenz um künstlich verknappte Mittel und
    Drittmittel-Hechelei. Mit diesem Engagement ist es bisher gelungen, der mit
    dem Dogma der "unternehmerischen Hochschule" forcierten Kultur des Misstrauens
    und der Konkurrenz eine neue Kultur der Solidarität entgegenzusetzen. So ist
    die Diskussion um die gesellschaftliche Verantwortung der Uni neu befeuert,
    was Ausdruck in den Protesten für eine Demokratisierung und bedarfsgemäße
    Grundfinanzierung der Wissenschaft und in der Studienreform findet.

    Davon will der Wissenschaftsrat nichts wissen. Ohne Offenlegung und Begründung
    für die Prämissen der Beurteilung werden in der Empfehlung als Maßstäbe die
    (internationale) Sichtbarkeit durch Profilbildung, das Drittmittel-Volumen,
    die Wirtschafsanbindung und die Arbeitsmarktrelevanz gesetzt. Diese Orientierung
    an der Verwertbarkeit ist aber umfassend gescheitert. Gerade in der aktuell
    tiefen politischen und sozialen Polarisierung in der Gesellschaft muss
    Wissenschaft aufklärerisch eingreifen und sich beimessen, politisch wirksam
    zu sein. Die Empfehlungen des Wissenschaftsrats sind eben durch das
    Verkennen der entwickelten Realität zutiefst anachronistisch.

    Das Studierendenparlament nimmt sich vor diesem Hintergrund erstmal die
    Empfehlungen des Wissenschaftsrats im Bereich "Studium und Lehre" vor. Denen
    entgegen engagiert sich das Studierendenparlament

    \begin{myitemize}
        \item für den Ausbau der Volluniversität, weil in der Vielfalt der Fächer
        die breite Wahrnehmung von historisch gewachsener und gesellschaftlicher
        Verantwortung zum Ausdruck kommt, und insbesondere die Kleinen Fächer in
        den Sprachwissenschaften Bedeutung für die diplomatische und zivile
        Verständigung der Menschen weltweit haben,
        \begin{myitemize}
            \item statt dem Wissenschaftsrat zu glauben, die Profilbildung sei
            wichtiger als Volluniversität, weil Verwertbarkeit und
            Wirtschafsanbindung wichtiger sei.
        \end{myitemize}
        \item für die Weiterentwicklung der Einheit von Forschung und Lehre mit
        der Orientierung an den gesellschaftlichen Herausforderungen wie
        Friedensentwicklung, Demokratisierung, Herstellung sozialer Gerechtigkeit
        sowie Inklusion aller Menschen,
        \begin{myitemize}
            \item statt wie der Wissenschaftsrat fordert, das Angebot der
            Sonderpädagogik zu überprüfen, da es zu geringe Personalausstattung
            und Studierendenzahl habe, und zu prüfen, ob Studiengänge in den
            Kleinen Fächern besser gestrichen werden sollten.
        \end{myitemize}
        \item für den Ausbau des Angebots "Deutsch als Fremdsprache", um die Uni
        Hamburg Studieninteressierten aus aller Welt stärker zu öffnen und die
        Kommiliton*innen durch die Möglichkeit des Spracherwerbs in die Lage zu
        versetzen, an dem politischen und kulturellen Leben an der Uni und in
        Hamburg teilzuhaben,
        \begin{myitemize}
            \item statt dem Wissenschaftsrat zu folgen, für höhere internationale
            Attraktivität die Unterrichtssprache (auf Englisch) umzustellen.
        \end{myitemize}
        \item für die weitere Belebung des internationalen Austauschs unter den
        Kommiliton*innen und Wissenschaftler*innen an der Uni Hamburg vor Ort,
        um die Kenntnisse in einer Fremdsprache und in anderen Kulturen auszubauen,
        \begin{myitemize}
            \item statt der Empfehlung des Wissenschaftsrat nachzukommen, in den
            Lehramtsstudiengängen ein verpflichtendes Auslandssemester einzuführen,
            besonders ohne jegliche Reflexion der prekären Lage der Studierenden.
        \end{myitemize}
        \item für den Ausbau von dauerhaften Personalstellen in der Lehre und
        damit die Verbesserung des Lehr-Lern-Verhältnisses, weil dadurch ein
        solidarisches und demokratisches Lernen gestärkt wird,
        \begin{myitemize}
            \item statt sich dem Wissenschaftsrat anzuschließen, ein gutes
            Betreuungsverhältnis aufgrund von niedrigen Studierendenzahlen zum
            Anlass zu nehmen, die Effizienz des Studiengangs in Frage zu stellen
            und zu fordern, Hochschulpaktmittel für Professuren und Lehrpersonal
            einzusetzen.
        \end{myitemize}
        \item für das Ende des Akkreditierung-Wahnsinns, von dem nur die privaten
        Akkreditierungsagenturen profitieren, und dem Ausbau der demokratischen
        Studiengangs-Organisation und -Evaluierung,
        \begin{myitemize}
            \item statt dem fast schon verzweifelten Fordern des Wissenschaftsrats
            nachzugeben, jetzt doch noch, auch nach dem
            Bundesverfassungsgerichtsurteil, die aktuellen Regelung zur
            Akkreditierung ginge nicht mit der im Grundgesetz verankerten Freiheit
            von Wissenschaft einher, die Studiengänge an der Uni Hamburg zu
            akkreditieren.
        \end{myitemize}
    \end{myitemize}

    Das Studierendenparlament nimmt die Empfehlung des Wissenschaftsrat zum Anlass,
    alle zu ermutigen, die konformitätserzwingenden Instrumente wie den
    Wissenschaftsrat zu verlachen, weil diese keine Perspektive haben die aktuell
    tiefe Krise human zu lösen, und sich selbst daran zu machen, mit Wissenschaft
    dazu beizutragen, für alle ein menschenwürdiges Leben zu realisieren. Die
    Möglichkeiten dafür sind aktuell mit dem erarbeiteten Reichtum so gut wie
    noch nie!

\end{document}
